% Technical basis: Apache name-based virtual hosting; MDN web server.
\subsubsection{প্রকাশের জন্য পাবলিক IP: কার ঠিকানা?}
অনলাইনে ওয়েবসাইট দেখতে ব্যবহারকারীর অনুরোধ পৌঁছানোর ব্যবস্থা থাকা প্রয়োজন। হোস্টিংয়ের সার্ভার সেই ব্যবস্থা দেয়; ডোমেইন ও DNS ঠিকানা খুঁজতে সাহায্য করে। শিক্ষার্থী হোস্টিং সেবা ব্যবহার করলে নিজের কম্পিউটারের জন্য আলাদা পাবলিক IP কিনে সেটি সবসময় চালু রাখতেই হবে এমন নয়।

\begin{keypointbox}{এক IP-তে একাধিক ওয়েবসাইট}
শেয়ারড হোস্টিংয়ে এক সার্ভারের একই IP ব্যবহার করে একাধিক ডোমেইনের সাইট চালানো যায়। অনুরোধে কোন ডোমেইন চাওয়া হয়েছে দেখে সার্ভার উপযুক্ত সাইটের ফাইল পাঠায়। তাই প্রতি ওয়েবসাইটের জন্য আলাদা IP বাধ্যতামূলক ধরে উত্তর লিখবে না। ডোমেইন নিবন্ধন করলেই ফাইল আপলোড বা হোস্টিং সম্পন্ন হয় না—প্রকাশের ধাপগুলো আলাদা।
\end{keypointbox}
